\section{Taxonomy of Dynamic Aspects in Robotic World Models}\label{sec:taxonomy_dynamics}
Before analyzing how factor graphs and scene graphs represent, construct and exploit non-stationary environments, we first set the vocabulary this review adopts to describe \textit{dynamics themselves}. 
The literature is not consistent here: the same underlying phenomenon, e.g., a chair being moved between visits, is variously described as a \textit{semi-static change} \cite{gomez_object-based_2020}, a \textit{movable} object \cite{brink_using_2016} or a node with a \textit{persistence} attribute\cite{behrens_lost_2025}.
Establishing a shared characterization upfront allows us to classify representations, pipelines, and applications along consistent axes for the remainder of the review.

We organize dynamics along three axes: 
\emph{when} a change occurs relative to observation (Temporal Horizon),
\emph{what} kind of entity undergoes it (Entity Typology),
and \emph{how} the change manifests in the representation (State Expression).
A fourth property, the \emph{predictability} of a dynamic entity's behavior, is briefly noted at the end of this section since it is only sparsely addressed by the surveyed literature.

\subsection{Primary Axis: Temporal Horizon}\label{sec:taxonomy_temporal_horizon}

\begin{table}[t]
\centering
\scriptsize
\setlength{\tabcolsep}{3pt}
\renewcommand{\arraystretch}{1.15}
\caption{Primary temporal-horizon axis used to organize dynamics in graph-based robot world models.}\label{tab:temporal_horizon_taxonomy}
    \begin{tabularx}{\columnwidth}{@{}p{0.15\columnwidth} p{0.12\columnwidth} X@{}}
    \toprule
    \textbf{Tier} & \textbf{Observation-relative criterion } & \textbf{Description} \\
    \midrule
    Short-Term Dynamics 
    &  Evolves during continuous observation 
    &  Motion must be estimated during observation, such as humans, vehicles, or robots \\ 
    \midrule
    Semi-Static Changes 
    &  Stable within observation, changes between observations/revisits 
    &  Relocation, appearance/disappearance, or persistent state transitions such as open/closed door \\ 
    \midrule
    Long-Term Structural Changes 
    &  Persistent structural modification across sessions 
    &  Changes to the built environment itself, such as layout changes, renovations, or seasonal structural variations  \\
    \bottomrule
    \end{tabularx}
\end{table}
% \begin{table}[t]
% \centering
% \caption{Primary temporal-horizon axis used to organize dynamics in graph-based robot world models.}\label{tab:temporal_horizon_taxonomy}
%     \begin{tabularx}{\columnwidth}{@{}p{0.12\columnwidth} p{0.12\columnwidth} X X@{}}
%     \toprule
%     \textbf{Tier} & \textbf{Timescale} & \textbf{Description} & \textbf{Required model mechanism} \\
%     \midrule
%     Dynamic Agents 
%     & Continuous / Now 
%     & Independently moving entities whose pose changes during observation, such as humans, vehicles, or robots 
%     & Real-time tracking, motion-aware data association, and often joint SLAMMOT-style estimation \\
%     \midrule
%     Semi-Static Changes 
%     & Minutes--Days 
%     & Entities that are static during a single observation but may be relocated, added, removed, or change state between observations 
%     & Persistence reasoning, object permanence, containment relations, and state-update policies \\
%     \midrule
%     Long-Term Structural Changes 
%     & Days--Years 
%     & Persistent modifications to the built environment itself, such as layout changes, renovations, or seasonal structural variations 
%     & Multi-session alignment, change detection, map repair, and lifelong graph maintenance \\
%     \bottomrule
%     \end{tabularx}
% \end{table}

The dimension that structures most consistently both individual models and this review is \textit{when} a change occurs relative to the robot's observation stream, since it directly governs the memory allocation, data association strategy, and update frequency a world model must support. 
We distinguish three tiers:

\begin{itemize}
    \item \textbf{Short-Term Dynamics} ( within-observation evolution): humans, moving vehicles, and other independently moving entities whose pose changes within the timescale of a single operating session, requiring real-time tracking and motion-aware data association \cite{bescos_dynaslam_2021} \cite{morris_dynosam_2025} \cite{schmid_khronos_2024}.
    \item \textbf{Semi-Static Changes} ( between-observation state changes): objects that are stationary during any single observation but relocated, added, removed, or state-transitioned between observations, e.g., furniture rearrangement or a door's open/closed state \cite{gomez_object-based_2020} \cite{giberna_dynemo-slam_2025} \cite{schmid_khronos_2024}.
    \item \textbf{Long-Term Structural Changes} ( persistent cross-session evolution): modifications to the built environment itself, e.g., renovations or seasonal layout shifts, requiring inter-session alignment and lifelong map maintenance \cite{walcott-bryant_dynamic_2012} \cite{kim2022lt}.
\end{itemize}
These tiers classify motion and state-change processes, rather than physical entities themselves.
Temporal Horizon is defined relative to the robot's observation process and world-model update cycle, rather than by a universal elapsed-time threshold.
\hl{Although these regimes are typically associated with different orders of magnitude in elapsed time, such temporal ranges are only indicative and do not define classification boundaries.
The classification depends on the frequency and continuity of observations, the operational context of the robot, and the persistence and structural nature of the environment change. 
For instance, the same object displacement may be perceived as continuous motion by a robot directly observing it, or a semi-static change by a robot revisiting the scene after the displacement occurred. }
A change process is classified as \textbf{short-term} when its state evolves during continuous observation and must therefore be estimated or tracked online; as \textbf{semi-static} when the entity is approximately stationary within an observation but its state differs between observations or revisits; and as \textbf{long-term structural} when the change constitutes a persistent modification of the environment's structural configuration across sessions or deployments. 
Consequently, different dynamic processes associated with the same physical entity may belong to different tiers, and may coexist within the same world model. 
For example, the articulated trajectory of a door while it is opening is a short-term dynamic process, the resulting open/closed state retained between observations is a semi-static state change, and installation, removal, or persistent reconfiguration of the door or doorway constitutes long-term structural evolution.
The taxonomy is therefore operational rather than object-specific. The three regimes imply different representation and maintenance requirements, ranging from continuous motion estimation and tracking, through state persistence and transition reasoning, to multi-session structural change detection and map restructuring.

% \hl{These tiers should be interpreted as operational temporal regimes rather than mutually exclusive classes of physical entities.
% The classification applies to a particular motion or state-change process relative to the robot's observation horizon and task requirements. Consequently, the boundaries between tiers are intentionally approximate and application-dependent.
% The same physical entity may therefore participate in multiple temporal regimes. 
% A door, for instance, can exhibit short-term articulated motion while being opened or closed, retain semi-static open/closed state across observations, and undergo a long-term structural change if the door or doorway is permanently installed, removed, or reconfigured. 
% Accordingly, a world model may represent different state variables or change process associated with the same entity at different temporal horizons. 
% }

This tiering is not merely descriptive: it maps directly onto the mechanism a world model must implement, from the tightly-coupled SLAMMOT formulations used for Short-Term Dynamics, through the persistence and containment reasoning used for Semi-Static Changes, to multi-session change detection used for Long-Term Structural Changes. 
We adopt this three-tier distinction as the primary organizing axis for the remainder of the paper.
A summary is depicted in Table \ref{tab:temporal_horizon_taxonomy}.

\subsection{Secondary Axes: Entity Typology and State Expression}\label{sec:taxonomy_secondary_axes}
Within each temporal tier, two further properties are useful for characterizing \textit{what} is changing and \textit{how} that change is expressed in the graph. 
Table \ref{tab:dynamics_taxonomy} summarizes both.

\textbf{Entity Typology} distinguishes dynamic elements by the structural assumption their motion model requires. 
\textit{Rigid} entities undergo $SE(3)$ pose change with constant geometry and dominate the surveyed factor-graph formulation \cite{bescos_dynaslam_2021} \cite{qiu2022airdos}.
\textit{Articulated} entities are constrained by internal kinematic chains, relevant both to manipulable objects such as doors and drawers \cite{zhang_open-vocabulary_2025} and to parametric human body models \cite{henning_bodyslam_2022} \cite{henning_bodyslam_2023}.
\textit{Deformable} entities, whose shape or topology changes non-rigidly, remain largely outside the scope of the probabilistic graph-based systems surveyed here.

\textbf{State Expression} distinguishes \textit{how} a change manifests once represented. 
\textit{Geometric-kinematic} changes are variations in position, velocity, or boundary and are the traditional target of factor graph estimation.
\textit{Semantic-affordance} changes affect an entity's functional or operational state, e.g., a light turning on or off, and are typically encoded as node attributes in 3D scene graphs \cite{catalano_3d_2025} \cite{olivastri_multi-modal_2024}.
\textit{Topological} changes modify the connectivity of the environment itself, e.g., a closed door removing the traversable edge between two rooms, and are of particular relevance to hierarchical scene graph maintenance \cite{catalano_3d_2025}.

\begin{table}[t]
\centering
\scriptsize
\setlength{\tabcolsep}{3pt}
\renewcommand{\arraystretch}{1.15}
\caption{Secondary axes used to characterize dynamic entities within each
temporal tier.}\label{tab:dynamics_taxonomy}
    \begin{tabularx}{\columnwidth}{@{}l l X@{}}
    \toprule
    \textbf{Axis} & \textbf{Category} & \textbf{Description} \\
    \midrule
    \multirow{3}{*}{\shortstack[l]{Entity\\Typology}}
     & Rigid       & $SE(3)$ pose change, constant geometry \\
     & Articulated & Kinematic chain, constrained internal DoF \\
     & Deformable  & Non-rigid shape/topology change \\
    \midrule
    \multirow{3}{*}{\shortstack[l]{State\\Expression}}
     & Geometric-Kinematic  & Position, velocity, boundary change \\
     & Semantic-Affordance  & Functional/operational state change \\
     & Topological          & Connectivity change in the scene graph \\
    \midrule
    \midrule
    \multirow{3}{*}{\shortstack[l]{Behavioral\\Intent}}
     & Deterministic  & Motion strictly governed by fixed laws \\
     & Physics-Based  & Motion governed by physics-based probabilities \\
     & Goal-Driven    & Motion governed by internal decisions or social compliance \\
    \bottomrule
    \end{tabularx}
\end{table}

\subsection{Adopted Characterization}\label{sec:taxonomy_adopted_characterization}
We adopt \textbf{Temporal Horizon} as the primary organizing dimension of this review, using the three-tier distinction to structure this paper's analysis.
The tiers are not mutually exclusive labels for entities or systems. They characterize change processes at different observation horizons, and a single entity or world model may simultaneously instantiate more than one tier.
Entity Typology and State Expression are used as secondary, descriptive dimensions throughout, applied where relevant to characterize individual systems without constituting an independent organizational axis.
A fourth property, \textit{Behavioral Intent}, that is the degree to which a dynamic entity's future motion is deterministic, physics-based or goal-driven, is addressed by only a small subset of the reviewed literature \cite{catalano_aion_2025} \cite{gorlo_long-term_2024} and remains largely decoupled from probabilistic world-model construction.
%We revisit it explicitly as an open direction in Section \ref{sec:open_challenges}

\begin{comment}

When building a general taxonomy for dynamic aspects in robotic world models, the goal is to categorize how a robot perceives, tracks, predicts, and reasons about a world that is not frozen in time. A general taxonomy of dynamic aspects in robotic world models can be structured around four interconnected pillars. This framework defines how a robot perceives, maintains, and projects environmental changes over time:
\begin{enumerate}
    \item \textbf{Entity Typology (The "What")} This dimension classifies the structural and material properties of the moving or changing elements, which dictates how they must be mathematically modeled:
    \begin{itemize}
        \item Rigid Dynamics: $SE(3)$ pose changes where geometry remains constant (e.g., vehicles, pedestrians).
        \item Kinematic Chains: Articulated objects with constrained internal degrees of freedom (e.g., opening doors, drawers).
        \item Deformable Elements: Non-rigid materials that change shape or topology (e.g., cables, cloth, liquids).
        \item Substance Fields: Continuous environmental phenomena lacking discrete boundaries (e.g., smoke, changing illumination).
    \end{itemize}
    \item \textbf{State Expression (The "How")} Dynamism manifests as changes in either the spatial layout or the internal characteristics of the world:
    \begin{itemize}
        \item Geometric-Kinematic Dynamics: Trackable variations in position, velocity, and physical boundaries.
        \item Semantic-Affordance Dynamics: Changes in the operational state or functional availability of an object (e.g., a stove turning "ON", a container becoming "Full").
        \item Topological Dynamics: Modifications to the connectivity of the environment (e.g., a closed door breaking the traversable edge between a hallway and a room). 
     \end{itemize}
    \item \textbf{Temporal Horizons (The "When")}
    This dimension governs the memory allocation, data association, and updating frequency required by the world model:
    \begin{itemize}
     \item Transient (Short-Term): High-frequency, immediate changes (seconds to minutes) requiring real-time tracking and collision avoidance.
     \item Episodic (Medium-Term/Semi-Static): Intermittent changes occurring between observations (hours to weeks), requiring change detection and graph-mutation (e.g., relocated furniture).
     \item Structural (Long-Term): Permanent or seasonal changes (months to years) requiring lifelong map maintenance and the discarding of stale environmental anchors.
    \end{itemize}
    \item \textbf{Behavioral Intent (The "Why")}
    This categorizes the predictability of the dynamic element, moving from reactive tracking to proactive anticipation:
    \begin{itemize}
    \item Deterministic: Motion strictly governed by fixed laws or periodic schedules (e.g., a conveyor belt).
    \item Stochastic-Kinematic: Motion governed by physics-based probabilities (e.g., a vehicle maintaining a constant velocity model with noise).
    \item Intentional/Social: Goal-driven motion governed by internal decisions or social compliance (e.g., a human walking toward a specific tool, obeying traffic rules).
    \end{itemize}
\end{enumerate}
\end{comment}
